\section{Comportamiento local de una función holomorfa}
\subsection{Comportamiento local de una función holomorfa}

Sea $\Omega \subset \mathbb{C}$ un dominio y sea $f \in \mathcal{H}(\Omega)$.\\

Nos interesa saber cuantos ceros tiene $f$ en un disco $D$ con $\overline{D} \subset \Omega$, o en general, en $\Int(\gamma)$ con $\gamma$ una curva cerrada homótopa a 0 en $\Omega$.\\

Para ello, nos va a ayudar la derivada logarítmica de $f$, es decir, la función:
\[
g(z) := \frac{f'(z)}{f(z)}
\]

Supongamos en primer lugar que $f$ tiene un numero finito de ceros en $\Omega$, siendo estos $\{ z_1, z_2, \ldots, z_n \}$, donde cada cero está repetido tantas veces como indica, sin multiplicidad.\\

Aplicando n veces la \cref{prop:ceros_aislados} tenemos que $\exists g \in \mathcal{H}(\Omega)$ tal que:
\[
\begin{cases}
    g(z_j) \neq 0 \quad \forall j = 1, \ldots, n \\
    f(z) = (z-z_1)(z-z_2) \cdots (z-z_n) g(z) \quad \forall z \in \Omega
\end{cases}
\]

Como $f$ no se anula en $\Omega \setminus \{ z_1, z_2, \ldots, z_n \}$, tenemos que $g(z) \neq 0$ en dicho conjunto.\\

Por tanto $f(z) = (z-z_1)(z-z_2) \cdots (z-z_n) g(z)$ con $g \in \mathcal{H}(\Omega)$ y $g(z) \neq 0$ en $\Omega$, entonces:
\[
f(z) =
\left( \prod_{j=1}^{n} (z-z_j) \right) g(z) \implies
f'(z) =
\left( \prod_{j=1}^{n} (z-z_j) \right) g'(z) +
\left( \sum_{k=1}^{n} \prod_{\substack{j=1 \\ j \neq k}}^{n} (z-z_j) \right) g(z)
\]
Así, tenemos que:
\[
\frac{f'(z)}{f(z)} =
\frac{g'(z)}{g(z)} + \sum_{j=1}^{n} \frac{1}{z-z_j} \quad \forall z \in \Omega \setminus \{ z_1, z_2, \ldots, z_n \}
\]

Sea ahora una curva $\gamma$ cerrada y homótopa a 0 en $\Omega$, que además no pasa por los ceros de $f$, es decir $z_j \notin \gamma^* \quad \forall j = 1, \ldots, n$, entonces tenemos que:
\[
\int_\gamma \frac{f'(z)}{f(z)} dz =
\int_\gamma \frac{g'(z)}{g(z)} dz + \sum_{j=1}^{n} \int_\gamma \frac{1}{z-z_j} dz
\]

Y ahora como $\frac{g'}{g} \in \mathcal{H}(\Omega)$, por ser $g \in \mathcal{H}(\Omega)$ y $g(z) \neq 0$ en $\Omega$, entonces aplicando el teorema de Cauchy para curvas homótopas, tenemos que:
\[
\int_\gamma \frac{g'(z)}{g(z)} dz = 0
\]
Con lo cual:
\[
\frac{1}{2 \pi i} \int_\gamma \frac{f'(z)}{f(z)} dz =
\sum_{j=1}^{n} \frac{1}{2 \pi i} \int_\gamma \frac{1}{z-z_j} dz =
\sum_{j=1}^{n} \Ind_\gamma (z_j)
\]

Si $\gamma$ es una curva simple positiva, entonces $\Ind_\gamma(z_j) = 1$ en el interior de $\gamma$ y 0 en el exterior, con lo cual:
\[
\frac{1}{2 \pi i} \int_\gamma \frac{f'(z)}{f(z)} dz =
\# \{ \text{ceros de } f \text{ en } \Int(\gamma) \}
\]

Si $f$ tiene infinitos ceros en $\Omega$, vamos a probar que ocurre lo mismo. Para ello necesitamos un lema (que usaremos también más adelante):

\begin{lema}
    Sea $\Omega \subset \mathbb{C}$ un dominio y sea $\gamma$ una curva cerrada homótopa a $0$ en $\Omega$. Entonces existe un dominio acotado $\Omega_1 \subset \Omega$ tal que $\overline{\Omega}_1 \subset \Omega$, $\gamma \sim 0$ en $\Omega_1$, y además $\Ind_\gamma(z) = 0$ para todo $z \notin \Omega_1$ (en particular, $\Int(\gamma) \subset \Omega_1$).
\end{lema}

\begin{proof}
    Dado que $\gamma \sim 0$ en $\Omega$, existe una aplicación continua 
    \[
        \varphi : [0,1]^2 \to \Omega
    \]
    que define la homotopía entre $\gamma$ y el lazo constante. Sea entonces 
    \[
        K = \varphi([0,1]^2),
    \]
    que es un conjunto compacto y conexo contenido en $\Omega$ (por ser imagen de un compacto por una aplicación continua). En particular, se cumple que $\Int(\gamma) \subset K$.

    Por el Lema 1 del Tema 8, existe $r > 0$ tal que
    \[
        \Omega_0 := \bigcup_{z \in K} D(z,r) \subset \Omega,
    \]
    Definimos ahora
    \[
        \Omega_1 := \bigcup_{z \in K} D(z, r/2) = K + D(0, r/2) \subset \Omega_0 = K + D(0,r) \subset \Omega.
    \]
    El conjunto $\Omega_1$ es abierto y conexo, pues $K$ lo es. Además,
    \[
        \overline{\Omega}_1 
        = \overline{K + D(0, r/2)} 
        = \overline{K} + \overline{D}(0, r/2)
        = K + \overline{D}(0, r/2)
        \subset K + D(0, r) 
        = \Omega_0 
        \subset \Omega.
    \]
    Por tanto, $\Omega_1$ es acotado, ya que es suma de dos conjuntos acotados ($K$ y $D(0, r/2)$).

    Como $\Int(\gamma) \subset K \subset \Omega_1$ y $\Omega_1 \subset \Omega$ es un dominio, se tiene que $\gamma \sim 0$ en $\Omega_1$. Finalmente, si $z \notin \Omega_1$, entonces $\Ind_\gamma(z) = 0$.
\end{proof}


Tenemos entonces el siguiente resultado (que es un caso particular del Principio del Argumento):

\begin{teorema}
    Sea $\Omega \subset \mathbb{C}$ un dominio y sea $f \in \mathcal{H}(\Omega)$, $f \not\equiv 0$. Sea $\{z_j : j \in J \}$ (con $J \subset \mathbb{N}$) el conjunto de ceros de $f$ en $\Omega$ repetidos cada uno de ellos tantas veces como indica su multiplicidad. Entonces para cada curva cerrada $\gamma$ homótopa a 0 en $\Omega$ que no pasa por los ceros de $f$, se cumple que:
    \[
        \frac{1}{2 \pi i} \int_\gamma \frac{f'(z)}{f(z)} dz =
        \sum_{j \in J} \Ind_\gamma (z_j)
    \]
\end{teorema}

\begin{proof}
    Sea $\gamma$ una curva que cumple las hipótesis del teorema. Sea $\Omega_1 \subset \Omega$ el dominio que asegura el lema anterior. Tenemos que, como $\overline{\Omega}_1$ es compacto y $Z_j$ no tiene punots de acumulación en $\Omega$, solamente hay una cantidad finita de ceros de $f$ en $\overline{\Omega}_1$.\\
    Puesto que $\gamma \sim 0$ en $\Omega_1$. Por lo probado anteriormente como $f \in \mathcal{H}(\Omega_1)$ y tiene solamente un número finito de ceros en $\Omega_1$, se cumple que:
    \[
        \frac{1}{2 \pi i} \int_\gamma \frac{f'(z)}{f(z)} dz =
        \sum_{z_j \in \Omega_1} \Ind_\gamma (z_j)
    \]
    Ahora bien, si $z_j \notin \Omega_1$, por el lema anterior se cumple que $\Ind_\gamma(z_j) = 0$, con lo cual:
    \[
        \sum_{\substack{j \in J \\ z_j \in \Omega_1}} \Ind_\gamma (z_j) = \sum_{j \in J} \Ind_\gamma(z_j)
    \]
\end{proof}

\begin{observación}
    \vspace{-2em}
    \begin{enumerate}
        \item En las hipótesis del teorema anterior, si tenemos que además $\gamma$ es una curva simple positiva, entonces:
        \[
            \frac{1}{2 \pi i} \int_\gamma \frac{f'(z)}{f(z)} dz =
            \# \{ \text{ceros de } f \text{ en } \Int(\gamma) \}
        \]
        \item 
        \[
            \frac{1}{2 \pi i} \int_\gamma \frac{f'(z)}{f(z)} dz = \frac{1}{2 \pi} \int_b^a \frac{f'(\gamma(t))}{f(\gamma(t))} \gamma'(t) dt =^*
        \]
        tomando $\gamma : [a,b] \to \mathbb{C}$ parametrización de $\gamma$. Ahora, haciendo la composición $h = f \circ \gamma : [a,b] \to \mathbb{C}$, tenemos que:
        \[
            =^* \frac{1}{2 \pi i} \int_a^b \frac{(f \circ \gamma)'(t)}{(f \circ \gamma)(t)} dt = \Ind_{f \circ \gamma}(0)
        \]
        Luego el teorema anterior nos dice que 
        \[
            \frac{1}{2 \pi i} \int_\gamma \frac{f'(z)}{f(z)} dz = \Ind_{f \circ \gamma}(0) = \sum_{j \in J} \Ind_\gamma (z_j)
        \]
        o sea, que $\frac{1}{2 \pi i} \int_\gamma \frac{f'(z)}{f(z)} dz$ cuenta cuántas veces rodea la curva $f \circ \gamma$ al $0$.
    \end{enumerate} \label{obs:comportamiento_local}
\end{observación}

\subsection{caracterizacion del comportamiento local de una funcion holomorfa}

\begin{teorema} [Comportamiento local de las funciones holomorfas]
    Sea $\Omega \subset \mathbb{C}$ un dominio y sea $f \in \mathcal{H}(\Omega)$, $f$ no constante. Sea $z_0 \in \Omega$ y llamamos $w_0 = f(z_0)$. Si $z_0$ es un cero de orden $k$ de la función $f(z) - w_0$, entonces existen $\delta, r >0$ tales que:
    \[
        \overline{D}(z_0,r) \subset \Omega, \quad \text{y} \quad D(w_0, \delta) \subset f(D(z_0,r))
    \]
    y para cada $w \in D(w_0,\delta)$ la función $f-w$ tiene exactamente $k$ ceros en $D(z_0,r)$ (que en el caso $w \neq w_0$ son todos distintos).
\end{teorema}

\begin{proof}
    Como $\Omega$ es abierto y $z_0 \in \Omega$, existe $r_1 > 0$ tal que $\overline{D}(z_0, r_1) \subset \Omega$.

    Dado que $f - w_0$ no es constante, el punto $z_0$ es un cero aislado de dicha función. Por tanto, existe $r_2 \in (0, r_1)$ tal que
    \[
        f(z) - w_0 \neq 0 \quad \text{para todo } z \in D(z_0, r_2) \setminus \{z_0\}.
    \]
    Como $f$ no es constante, la derivada $f'$ no puede ser idénticamente cero en $D(z_0, r_2)$. Luego $Z_{f'}$ no tiene puntos de acumulación en $D(z_0, r_2)$; es más: tiene una cantidad finita de puntos (pues $D(z_0,r_2)$ está acotado), por lo que podemos elegir $r_3 \in (0, r_2)$ tal que
    \[
        f'(z) \neq 0, \quad \forall z \in D(z_0, r_3) \setminus \{z_0\}.
    \]
    Fijamos $r \in (0, r_3)$ y consideramos la circunferencia $\gamma = C(z_0, r)$, orientada positivamente.  
    Tenemos que $w_0 \notin f(C(z_0, r))$ al ser $f(z) - w_0 \neq 0$ en $D(z_0, r_2) \setminus \{z_0\} \supset C(z_0, r)$.\\
    Como $f(C(z_0, r))$ es compacto (imagen continua de un compacto), entonces $\mathbb{C} \setminus f(C(z_0, r))$ es abierto y contiene a $w_0$. Por tanto, existe $\delta > 0$ tal que
    \[
        D(w_0, \delta) \subset \mathbb{C} \setminus f(C(z_0, r)).
    \]
    o que es lo mismo,
    \[
        D(w_0, \delta) \cap f(C(z_0, r)) = \emptyset.
    \]
    Esto significa que el disco $D(w_0, \delta)$ queda completamente dentro de una de las componentes conexas del complemento de $f(C(z_0, r))$. Denotando esa componente por $U$, se tiene
    \[
        D(w_0, \delta) \subset U \subset \mathbb{C} \setminus f(C(z_0, r)).
    \]

    Como el índice $\Ind_{f \circ \gamma}(w)$ es constante en cada componente conexa del complemento de la imagen de la curva, se sigue que
    \[
        \Ind_{f \circ \gamma}(w) = \Ind_{f \circ \gamma}(w_0), \quad \forall\, w \in D(w_0, \delta).
    \]

    En particular, aplicando la definición de la función índice,
    \[
        \Ind_{f \circ \gamma}(w_0)
        = \frac{1}{2\pi i} \int_{f \circ \gamma} \frac{1}{z - w_0} \, dz
        = \frac{1}{2\pi i} \int_\gamma \frac{f'(z)}{f(z) - w_0} \, dz
        = \frac{1}{2\pi i} \int_\gamma \frac{(f - w_0)'(z)}{(f - w_0)(z)} \, dz.
    \]

    Por el apartado (1) de la \cref{obs:comportamiento_local}, si $\gamma$ es una curva simple orientada positivamente, esta integral cuenta el número de ceros de $f - w_0$ en el interior de $\gamma$. Así,
    \[
        \Ind_{f \circ \gamma}(w_0)
        = \# \{ \text{ceros de } f - w_0 \text{ en } D(z_0, r) \}
        = k.
    \]

    De manera análoga, para cualquier $w \in D(w_0, \delta)$ tenemos
    \[
        \Ind_{f \circ \gamma}(w)
        = \frac{1}{2\pi i} \int_\gamma \frac{f'(z)}{f(z) - w} \, dz
        = \frac{1}{2\pi i} \int_\gamma \frac{(f - w)'(z)}{(f - w)(z)} \, dz
        = \# \{ \text{ceros de } f - w \text{ en } D(z_0, r) \}.
    \]
    Como el índice es constante en $D(w_0, \delta)$, obtenemos
    \[
        \# \{ \text{ceros de } f - w \text{ en } D(z_0, r) \} = k, \quad \forall w \in D(w_0, \delta).
    \]

    Finalmente, si $w \neq w_0$, estos ceros son todos simples (y por tanto distintos), porque $f'(z) \neq 0$ en $D(z_0, r) \setminus \{z_0\}$, lo que impide que se produzcan ceros múltiples.
\end{proof}

Consecuencias importantes del teorema anterior son los siguientes resultados:

\subsection{Teorema de la aplicación abierta}

\begin{teorema} [Teorema de la aplicación abierta]
    Sea $\Omega \subset \mathbb{C}$ un dominio y sea $f \in \mathcal{H}(\Omega)$, $f$ no constante en $\Omega$. Entonces $f$ es una aplicación abierta, es decir, la imagen por $f$ de cualquier conjunto abierto de $\Omega$ es un conjunto abierto de $\mathbb{C}$.
\end{teorema}

\begin{proof}
    Veamos primero que $f(\Omega)$ es abierto. Sea $w_0 \in f(\Omega)$, entonces $\exists z_0 \in \Omega$ tal que $f(z_0) = w_0$. Por el teorema anterior, existen $\delta, r > 0$ tales que:
    \[
        \overline{D}(z_0,r) \subset \Omega, \quad \text{y} \quad D(w_0, \delta) \subset f(D(z_0,r)) \subset f(\Omega) \subset f(\Omega).
    \]
    Luego $f(z_0)$ es un punto interior de $\Omega$.\\
    Si $G \subset \Omega$ es abierto y conexo, entonces como $f \in \mathcal{H}(\Omega)$, $f$ es continua en $\Omega$, luego $f(G)$ es conexo. Además, $f$ es no constante en $\Omega$, luego $f$ es no constante en $G$ por el principio de identidad.\\
    Aplicando la primera parte concluimos que $f(G)$ es abierto.\\
    Por último, si $G \subset \Omega$ es abierto (no necesariamente conexo), entonces $G$ es unión de sus componentes conexas, que son todas abiertas, luego por la segunda parte de la demostración, la imagen de cada una de ellas es un abierto con o cual la unión de todas las imgágenes de las componentes conexas de $G$ (que es $f(G)$) es un abierto.
\end{proof}

Una consecuencia inmediata del teorema anterior es:

\begin{corolario}
    Sea $\Omega \subset \mathbb{C}$ un dominio y sea $f : \Omega \to \mathbb{C}$ una función holomorfa inyectiva en $\Omega$. Entonces $f : \Omega \to f(\Omega)$ es biyectiva y su inversa en continua en $f(\Omega)$.
\end{corolario}

\subsection{Teorema de la función inversa}

\begin{teorema} [Teorema de la función inversa]
    Sea $\Omega \subset \mathbb{C}$ un dominio y sea $f \in \mathcal{H}(\Omega)$. Sea $z_0 \in \Omega$ tal que $f'(z_0) \neq 0$. Entonces existen un entorno abierto conexo $U \subset \Omega$ de $z_0$ tal que:
    \begin{enumerate}
        \item $f$ es inyectiva en $U$.
        \item $W := f(U)$ es un abierto.
        \item La función $f |_U : U \to W$ tiene inversa $f^{-1} : W \to U$ que es holomorfa en $W$.
    \end{enumerate}
\end{teorema}

\begin{proof}
    Llamamos $w_0 = f(z_0)$. Como $f'(z_0) = (f-w_0)'(z_0) \neq 0$, entonces $z_0$ es un cero de orden 1 de la función $f - w_0$. Por el teorema del comportamiento local de las funciones holomorfas, existen $\delta, r > 0$ tales que:
    \[
        \overline{D}(z_0,r) \subset \Omega, \quad \text{y} \quad D(w_0, \delta) \subset f(D(z_0,r))
    \]
    además para cada $w \in D(w_0, \delta)$ la función $f-w$ tiene exactamente un cero en $D(z_0,r)$, es decir, $\exists! z \in D(z_0,r) : f(z) = w$ (lo cual implica que $f$ es inyectiva en $D(z_0,r)$).\\
    Por lo tanto, tomando:
    \[
        U := D(z_0,r) \cap f^{-1}(D(w_0, \delta))
    \]
    Asi, $U$ es un entorno abierto conexo de $z_0$ en $\Omega$ en el que $f$ es inyectiva.
    Como $U \subset \Omega$ es abierto, por el teorema de la aplicación abierta, $W := f(U)$ es un abierto.\\
    Tenemos
    \[
        f |_U : U \to W
    \]
    es holomorfa y biyectiva, luego su inversa $g = f^{-1} : W \to U$ es continua por el corolario anterior.\\
    Falta comprobar que $g$ es holomorfa en $W$. Sea $w_1 \in W = f(U)$, entonces $\exists z_1 \in U$ tal que $f(z_1) = w_1$. \\
    Tomando la sucesión $(w_n)_{n \in \mathbb{N}} \subset W \setminus \{ w_1 \}$ que converge a $w_1$, tenemos que, para cada $n \in \mathbb{N}$, existe un único $z_n \in U$ tal que $f(z_n) = w_n$.\\
    Como $w_n \to w_1$, y $g$ es continua en $w_1$, tenemos que $z_n = g(w_n) \to g(w_1) = z_1$.\\
    Luego $(z_n)_{n \in \mathbb{N}} \subset U \setminus \{ z_1 \}$ es una sucesión que converge a $z_1$.\\
    Por lo tanto,
    \[
        \frac{f(z_n) - f(z_1)}{z_n - z_1} = \xrightarrow{n \to \infty} f'(z_1) \neq 0
    \]
    Con lo cual
    \[
        \frac{g(w_n) - g(w_1)}{w_n - w_1} = \frac{z_n - z_1}{f(z_n) - f(z_1)} \underbrace{=}_{z_n \neq z_1} = \frac{1}{\frac{f(z_n)-f(z_1)}{z_n-z_1}} \xrightarrow{n \to \infty} \frac{1}{f'(z_1)} 
    \]
    Para que $f'(z_1) \neq 0$, entonces al elegir $U$ se cumpla que $f'(z) \neq 0$ para todo $z \in U$, cosa que podemos conseguir puesto que $f'$ es continua y $f'(z_0) \neq 0$.\\
\end{proof}

\begin{corolario}
    Sea $\Omega \subset \mathbb{C}$ un dominio y sea $f \in \mathcal{H}(\Omega)$. Si $f$ es inyectiva en $\Omega$, entonces $f' (z) \neq 0$ para todo $z \in \Omega$, y $f : \Omega \to f(\Omega)$ es biholomorfa (es decir, biyectiva, holomorfa y su inversa es holomorfa).
\end{corolario}

\begin{proof}
    Si $\exists z_0 \in \Omega$ tal que $f'(z_0) = 0$, entonces $z_0$ no es un cero de orden $k \geq 2$ de la función $f - w_0$ (con $w_0 = f(z_0)$). Entonces por el teorema del comportamiento local de las funciones holomorfas, existen $\delta, r > 0$ tales que para cada $w_1 \in D(w_0, \delta) \setminus \{ w_0 \}$ existen al menos dos puntos distintos $z_1, z_2$ tal que $f(z_1) = f(z_2) = w_1$, con lo cual $f$ no es inyectiva en $D(z_0,r) \subset \Omega$, contradicción.\\
    Por lo tanto, $f'(z) \neq 0$ para todo $z \in \Omega$. Luego por el teorema de la función inversa, $f$ es localmente invertible en cada punto de $\Omega$. Pero como $f$ es biyectiva de $\Omega$ en $f(\Omega)$, entonces $f$ es globalmente invertible, y su inversa es holomorfa en $f(\Omega)$. Luego $f : \Omega \to f(\Omega)$ es biholomorfa.
\end{proof}

Algo interesante es que en $\mathbb{C}$ si $f \in \mathcal{H}(\Omega)$:
\[
    f \text{ inyectiva en } \Omega \implies f' (z) \neq 0, \quad \forall z \in \Omega
\]

Mientras que si $f$ es derivable en $I \subset \mathbb{R}$, es al reves:
\[
    f' (x) \neq 0, \quad \forall x \in I \implies f \text{ inyectiva en } I
\]

Si $f \in \mathcal{H}(\Omega)$, con $\Omega \subset \mathbb{C}$ dominio, y $f$ no constante, podemos decir que localmente "$f$ es como $z^k$", ya que si $z_0 \in \Omega$ y $w_0 = f(z_0)$, entonces $z_0$ es un cero aislado de $f - w_0$ de orden $k \geq 1$, y por el teorema anterior $\exists g \in \mathcal{H}(\Omega)$ con $g(z_0) \neq 0$ y:
\[
f(z) - w_0 = (z - z_0)^k g(z)
\]
Como $g$ esta contenida en $\Omega$ y $g(z_0) \neq 0$, entonces $\exists r > 0$ tal que $g(z) \neq 0$ en $D(z_0,r)$, con lo cual:
\[
\exists h \in \mathcal{H}(D(z_0,r)) : (h(z))^k = g(z) \quad \forall z \in D(z_0,r)
\]
Y como $g$ no se anula en $D(z_0,r)$, que es simplemente conexo, $h$ tampoco lo hara en $D(z_0,r)$. Por tanto:
\[
f(z) - w_0 =
(z - z_0)^k g(z) =
(z - z_0)^k (h(z))^k =
\left((z - z_0) h(z)\right)^k \quad \forall z \in D(z_0,r)
\]

Si llamamos $\varphi(z) = (z - z_0) h(z) \forall z \in D(z_0,r)$, tenemos que $\varphi \in \mathcal{H}(D(z_0,r))$, $\varphi'(z) = h(z) + (z - z_0) h'(z) \forall z \in D(z_0,r)$, y por tanto $\varphi'(z_0) = h(z_0) \neq 0$. Luego, por el teorema de la función inversa, $\exists U$ entorno abierto de $z_0$ tal que:
\[
\varphi_{|U} : u \mapsto \varphi(u)
\]
Es una función biholomorfa, y por tanto homeomorfismo, obteniendo finalmente que:
\[
f(z) = (\varphi(z))^k + w_0 \quad \forall z \in U
\]

\begin{teorema}
    Sea $\Omega \subset \mathbb{C}$ un dominio y sea $(f_n)_{n \in \mathbb{N}}$ una sucesión de funciones en $\mathcal{H}(\Omega)$ que converge uniformemente en los compactos de $\Omega$ a una función $f \in \mathcal{H}(\Omega)$. Si cada $f_n$ es inyectiva en $\Omega$, entonces $f$ es o bien constante o bien inyectiva en $\Omega$.
\end{teorema}

\begin{proof}
    Supongamos que $f$ no es constante en $\Omega$. Veamos que $f$ es inyectiva en $\Omega$.\\
    Sea $z_0 \in \Omega$, queremos probar que $f(z) \neq f(z_0) = w_0 \quad \forall z \in \Omega \setminus \{ z_0 \}$.\\
    Sea $z_1 \in \Omega \setminus \{ z_0 \} \implies \exists r > 0$ tal que (al ser $\Omega \setminus \{ z_0 \}$ abierto) $\overline{D}(z_1,r) \subset \Omega \setminus \{ z_0 \}$.\\
    Ademas, $\overline{D}(z_1,r) \cap Z_{f-w_0}$ es finito, pues el primer conjunto es compacto y el segundo no tiene puntos de acumulación en $\Omega$.\\
    Ahora tomemos $0 < r_1 < r$ tal que $f - w_0$ no se anule en $\overline{D}(z_1,r_1) \setminus \{ z_1 \}$.\\
    Así, $f - w_0$ no se anula en $C(z_1,r_1)$, y por tanto:
    \[
    \min_{z \in C(z_1,r_1)} |f(z) - w_0| = \rho > 0
    \]
    Como $f_n \to f$ uniformemente en $C(z_1,r_1)$, y tambien $f_n - f_n(z_0) \to f - w_0$ uniformemente en $C(z_1,r_1)$, entonces:
    \[
    \exists N \in \mathbb{N} :
    \max_{z \in C(z_1,r_1)} |(f_n - f_n(z_0)) - (f - w_0)| > \frac{\rho}{2} \quad \forall n \geq N
    \]
    Sea ahora $\gamma = C(z_1,r_1)$ orientada positivamente, entonces el Nº de ceros de $f - w_0$ en $\Int(\gamma) = \mathring{C}(z_1,r_1) = D(z_1,r_1)$ es, segun la proposicion anterior:
    \[
    = \frac{1}{2 \pi i} \int_\gamma \frac{(f'(z)}{f(z) - w_0} dz \underbrace{=}_{\text{Tª Weierstrass}}
    \frac{1}{2 \pi i} \int_\gamma \lim_{n \to \infty} \frac{f_n'(z)}{f_n(z) - f_n(z_0)} dz \underbrace{=}_{\text{Convergencia uniforme en } C(z_1,r_1) \text{ (ejercicio)}}
    \]
    Usando que las derivadas tambien convergen uniformemente en los compactos:
    \[
    = \lim_{n \to \infty} \frac{1}{2 \pi i} \int_\gamma \frac{f_n'(z)}{f_n(z) - f_n(z_0)} dz =
    \lim_{n \to \infty} \# \{ \text{ceros de } f_n - f_n(z_0) \text{ en } D(z_1,r_1) \} = 0
    \]
    Siendo eso ultimo ya que $f_n$ es inyectiva en $\Omega$, y como $z_0 \notin D(z_1,r_1)$, entonces $f_n(z) \neq f_n(z_0) \quad \forall z \in D(z_1,r_1)$.\\
    Luego $f - w_0$ no tiene ceros en $D(z_1,r_1)$, y por tanto $f(z_1) \neq w_0$.\\
    Ahora por el Tª de Weierstrass, $f'_n \to f'$ uniformemente en los compactos de $\Omega$, en particular en $C=C(z_1,r_2)$, y se comprueba que:
    \[
    \frac{1}{f_n - f_n(z_0)} \to \frac{1}{f - w_0} \quad \text{uniformemente en } C
    \]
    Con lo cual:
    \[
    \frac{f_n'}{f_n - f_n(z_0)} \to \frac{f'}{f - w_0} \quad \text{uniformemente en } C
    \]
\end{proof}